\section{Details on Fair Division Settings}
\label{sec:settings-extra}

\subsection{Valuation Function Type}

A function $u: 2^M \to \mathbb{R}$ is
\begin{tightenum}
\item \emph{additive} if for any two disjoint sets $S, T \subseteq M$, we have $u(S \cup T) = u(S) + u(T)$.
    Equivalently, for every set $S \subseteq M$, we have $u(S) = \sum_{j \in S} u(\{j\})$.
\item \emph{subadditive} if for any two disjoint sets $S, T \subseteq M$, we have $u(S \cup T) \le u(S) + u(T)$.
\item \emph{superadditive} if for any two disjoint sets $S, T \subseteq M$, we have $u(S \cup T) \ge u(S) + u(T)$.
\item \emph{submodular} if for any $S, T \subseteq M$, we have $u(S \cup T) + u(S \cap T) \le u(S) + u(T)$.
\item \emph{supermodular} if for any $S, T \subseteq M$, we have $u(S \cup T) + u(S \cap T) \ge u(S) + u(T)$.
\item \emph{cancelable} if for any $T \subseteq M$ and $S_1, S_2 \subseteq M \setminus T$,
    we have $u(S_1 \cup T) > u(S_2 \cup T) \implies u(S_1) > u(S_2)$.
\item \emph{unit-demand} if $u(\emptyset) = 0$, and for any $\emptyset \neq S \subseteq M$,
    we have $u(S) \defeq \max_{j \in S} u(\{j\})$.
\end{tightenum}

Note that when $|M|=1$, $u$ belongs to all of these classes simultaneously.

\subsection{Marginal Values}

\begin{tightenum}
\item \emph{goods}: $v_i(j \mid S) \ge 0$ for all $S \subseteq M$, $j \in M \setminus S$, and $i \in N$.
\item \emph{chores}: $v_i(j \mid S) \le 0$ for all $S \subseteq M$, $j \in M \setminus S$, and $i \in N$.
\item \emph{positive}: $v_i(j \mid S) > 0$ for all $S \subseteq M$, $j \in M \setminus S$, and $i \in N$.
\item \emph{negative}: $v_i(j \mid S) < 0$ for all $S \subseteq M$, $j \in M \setminus S$, and $i \in N$.
\item \emph{bivalued}: There exist constants $a, b \in \mathbb{R}$ such that
    $v_i(j \mid S) \in \{a, b\}$ for all $S \subseteq M$, $j \in M \setminus S$, and $i \in N$.
\item \emph{binary}: $v_i(j \mid S) \in \{0, 1\}$ for all $S \subseteq M$, $j \in M \setminus S$, and $i \in N$.
\item \emph{negative binary}: $v_i(j \mid S) \in \{0, -1\}$ for all $S \subseteq M$, $j \in M \setminus S$, and $i \in N$.
\end{tightenum}

We can break up the class of bivalued instances into positive bivalued, negative bivalued,
binary, negative binary, and mixed bivalued
(mixed means that exactly one of $a$ and $b$ is positive and the other is negative).
